\documentclass[10pt]{article}
\usepackage[margin=1.8cm]{geometry}
\usepackage{graphicx}
\usepackage{amsmath}
\usepackage{booktabs}
\usepackage[hidelinks]{hyperref}
\setlength{\parskip}{2pt}\setlength{\parindent}{0pt}

\title{\vspace{-1.4cm}Neuron mechanisms IV: amplifiers pass parities through the linear
window of $\mathrm{relu}(\tanh)$\\
\large \texttt{wd3\_k0\_N2048}, step 500k --- post-nonlinearity units, uniform evaluation data}
\author{}\date{\vspace{-1.2cm}}

\begin{document}\maketitle\vspace{-1.9cm}

\section*{The puzzle and the measurement}
\emph{Amplifier} neurons (blocks 2--3) receive an already-assembled parity: at least $90\%$
of their pre-activation variance lies on their own $\chi_{d8}$ (276 active neurons) or
$\chi_{d4}$ (22), and they re-emit that same mode at high purity. Since the parity is a
$\pm1$ signal, one would expect $\phi=\mathrm{relu}(\tanh)$ to cut the negative half and
distort the feature. It does not, because the input is not centered. Writing the
pre-activation on the own parity $\chi$ as $z = c + \gamma\chi$ and measuring both constants
per neuron gives Table~\ref{tab} and Fig.~\ref{fig:amp}: \textbf{the DC offset exceeds the
mode amplitude}, $c > |\gamma|$, for 297 of 298 amplifiers (median $c/|\gamma|\approx1.4$).
Both conditional levels $c\pm|\gamma|$ are therefore positive --- $z$ never reaches the relu
corner. Moreover the levels are small: the 95th-percentile upper level is $0.11$ ($d8$) and
$0.22$ ($d4$), where $\tanh$ deviates from the identity by $z^2/3\leq0.4\%/1.6\%$. So on its
operating window the neuron is a linear wire, $a\approx c+\gamma\chi$.

\begin{table}[h!]\centering\small
\begin{tabular}{lcccccccc}
\toprule
 & $n$ & med.\ $c$ & med.\ $|\gamma|$ & $c>|\gamma|$ & $z$ 95\% & $\dfrac{\Delta a}{\Delta z}$ &
 $\dfrac{|\langle a\chi\rangle|}{|\gamma|}$ & vs.\ clipped $\tfrac{\tanh\gamma}{2}$ \\
\midrule
$d8$ amplifiers & 276 & 0.040 & 0.028 & 275/276 & 0.106 & 0.996 & 0.995 & $1.99\times$ \\
$d4$ amplifiers & 22 & 0.086 & 0.061 & 22/22 & 0.218 & 0.991 & 0.992 & $1.99\times$ \\
\bottomrule
\end{tabular}
\caption{Per-neuron fits of $z=c+\gamma\chi$ and linearity checks. $\Delta a/\Delta z$: output
swing over input swing between the two parity classes; $|\langle a\chi\rangle|/|\gamma|$:
output over input Walsh coefficient. Both are $\approx\!1$, i.e.\ twice what half-wave
rectification of a centered input would leave.}
\label{tab}
\end{table}

\begin{figure}[h!]\centering
\includegraphics[width=0.93\textwidth]{fig8_amplifier_linear.pdf}\vspace{-4pt}
\caption{Left: exemplar $d8$ amplifier --- the two conditional distributions of $z$ (for
$\chi=\pm1$) both sit on the positive segment of $\mathrm{relu}(\tanh)$, where the curve is
indistinguishable from the identity. Middle: population scatter of offset $c$ vs.\ amplitude
$|\gamma|$; every neuron but one lies above the $c=|\gamma|$ line, outside the clipped
region. Right: output vs.\ input coefficient --- the points follow the identity
(linear pass-through), not the half-wave prediction $\tanh(\gamma)/2$ (dashed).}
\label{fig:amp}
\end{figure}

\section*{Why the offset is there, and what ``amplification'' means}
The offset is not a learned bias (the block linear layers have none). It is inherited from
the previous rectification: upstream carriers of a parity are one-sided units with
activations $\approx t(1+\chi)/2$, i.e.\ each delivers the mode with a DC component of the
same size attached. Any weighted sum of such features arrives with $c\sim|\gamma|$
automatically --- the DC produced by one layer's relu is exactly what shields the feature
from the next layer's relu.

The gain is therefore not pointwise ($|\langle a\chi\rangle| = |\gamma|$, nothing is
amplified inside the neuron). The amplification is collective: the $V$-row coherently sums
many weak copies of $\chi$ spread across the 2048 residual channels into a single clean
$\gamma$ (reading the stream is non-destructive), and the neuron writes a fresh copy back at
purity $0.94$/$0.92$ (median share of $\mathrm{Var}\,a$ on the own mode) --- versus $0.044$
for the raw-bit gates that created $d8$. Hundreds of such units re-emitting in parallel grow
the total mode power of the stream from block to block. Both polarities are present
($148{+}/128{-}$ for $d8$, $8{+}/14{-}$ for $d4$), so downstream sums can recover the full
$\pm1$ swing and cancel the accumulated DC. This completes the contrast with the creator
mechanisms: gates and corner units sit \emph{on} the relu corner and use it to compute new
modes, while amplifiers are biased \emph{past} the corner so as to transport existing modes
unchanged.

\vspace{2pt}\footnotesize
Artifacts: \texttt{gdrive:parity/analysis\_output/neuron\_mechanisms/wd3\_k0\_N2048/step\_00500000/}
(\texttt{amplifier\_stats.csv}, this figure and report; companion $d2$/$d4$/$d8$/$d16$
reports in the same folder).
\end{document}
